\clearpage
\section{삼차방정식과 Cardano 공식의 군론적 구조}
\label{sec:cardano-galois-structure}

\begin{learninggoal}
일반 삼차식을 우울형 삼차식으로 바꾸고, $S_3\trianglerighteq A_3\trianglerighteq1$의 정규열이 판별식의 제곱근과 세제곱근의 첨가로 번역되는 과정을 이해한다. Lagrange 분해식을 계산하여 Cardano 공식을 유도한다.
\end{learninggoal}

$\charac K\ne2,3$이라 하자. 일반 일계수 삼차식
\[
  X^3+aX^2+bX+c
\]
에
\[
  X=Y-\frac a3
\]
를 대입하면
\[
  Y^3+pY+q
\]
꼴의 \emph{우울형 삼차식}(depressed cubic)을 얻는다. 따라서
\[
  f(X)=X^3+pX+q
\]
를 푸는 것으로 충분하다.

\subsection{정규열이 지시하는 두 단계}

$f$의 분해체를 $L$, 세 근을 $x_1,x_2,x_3$이라 하자. 일반적인 경우 갈루아군은 $S_3$이고 정규열은
\[
  S_3\trianglerighteq A_3\trianglerighteq1
\]
이다. 갈루아 대응에서 첫 단계는
\[
  K\subseteq L^{A_3}
\]
인 이차확장이고, 둘째 단계는
\[
  L^{A_3}\subseteq L
\]
인 순환 삼차확장이다.

삼차식의 판별식은
\[
  \Delta=-4p^3-27q^2
\]
이고
\[
  \delta=(x_1-x_2)(x_1-x_3)(x_2-x_3),
  \qquad
  \delta^2=\Delta.
\]
홀순열은 $\delta$의 부호를 바꾸고 짝순열은 고정하므로
\[
  L^{A_3}=K(\delta)
\]
이다. 따라서 첫 단계는 $\sqrt\Delta$의 첨가이다.

두 번째 단계는 순환 삼차확장이다. 이를 세제곱근 하나로 표현하려면 원시 세제곱근
\[
  \zeta^2+\zeta+1=0,
  \qquad \zeta\ne1
\]
을 준비한다. $\zeta\notin K$이면 먼저 $K(\zeta)$로 기준체를 확대한다.

\subsection{Lagrange 분해식}

\begin{definition}[삼차 Lagrange 분해식]
근의 순서를 하나 택하고
\[
  R=x_1+\zeta x_2+\zeta^2x_3,
\]
\[
  S=x_1+\zeta^2x_2+\zeta x_3
\]
라 둔다. 이를 삼차식의 \emph{Lagrange 분해식}(Lagrange resolvents)이라 한다.
\end{definition}

$3$-순환 $\sigma=(123)$에 대해
\[
  \sigma(R)=\zeta^2R,
  \qquad
  \sigma(S)=\zeta S.
\]
따라서
\[
  R^3,S^3\in K(\delta,\zeta).
\]
이것이 순환 삼차확장을 세제곱근으로 푸는 핵심이다.

\begin{proposition}[분해식의 곱과 세제곱]
\label{prop:cubic-resolvent-identities}
다음 항등식이 성립한다.
\[
  RS=-3p,
\]
\[
  R^3=-\frac{27}{2}q+\frac32\sqrt{-3}\,\delta,
\]
\[
  S^3=-\frac{27}{2}q-\frac32\sqrt{-3}\,\delta.
\]
동치로
\[
  \left(\frac R3\right)^3
  =-\frac q2+\sqrt{\left(\frac q2\right)^2
  +\left(\frac p3\right)^3},
\]
\[
  \left(\frac S3\right)^3
  =-\frac q2-\sqrt{\left(\frac q2\right)^2
  +\left(\frac p3\right)^3}.
\]
\end{proposition}

\begin{proof}
$\zeta+\zeta^2=-1$과 Vieta 관계
\[
  x_1+x_2+x_3=0,
  \qquad
  x_1x_2+x_1x_3+x_2x_3=p,
  \qquad
  x_1x_2x_3=-q
\]
를 사용한다. 직접 곱하면
\[
\begin{aligned}
RS
&=x_1^2+x_2^2+x_3^2
 +(\zeta+\zeta^2)(x_1x_2+x_1x_3+x_2x_3)\\
&=(x_1+x_2+x_3)^2
 -3(x_1x_2+x_1x_3+x_2x_3)\\
&=-3p.
\end{aligned}
\]

$R^3$을 전개하고 대칭부분과 교대부분을 분리하면
\[
  R^3=-\frac{27}{2}q+\frac32(\zeta-\zeta^2)\delta.
\]
원시 세제곱근에 대해
\[
  \zeta-\zeta^2=\sqrt{-3}
\]
인 제곱근의 선택을 할 수 있으므로 첫 식을 얻는다. $S^3$은 $\zeta$와 $\zeta^2$를 바꾸어 얻는다.

마지막으로
\[
  \Delta=-4p^3-27q^2
\]
이므로
\[
  \left(\frac32\sqrt{-3}\,\delta\right)^2
  =-\frac{27}{4}\Delta
  =27^2\left(
  \left(\frac q2\right)^2+\left(\frac p3\right)^3
  \right).
\]
이를 $27$로 나누면 표시된 동치식이 나온다.
\end{proof}

\subsection{역 Fourier 변환}

세 근은 $R,S$로부터 복원된다. 실제로
\[
\begin{pmatrix}
0\\R\\S
\end{pmatrix}
=
\begin{pmatrix}
1&1&1\\
1&\zeta&\zeta^2\\
1&\zeta^2&\zeta
\end{pmatrix}
\begin{pmatrix}
x_1\\x_2\\x_3
\end{pmatrix}
\]
이고 이 Fourier 행렬의 역행렬을 쓰면
\[
  x_1=\frac{R+S}{3},
\]
\[
  x_2=\frac{\zeta^2R+\zeta S}{3},
\]
\[
  x_3=\frac{\zeta R+\zeta^2S}{3}.
\]

\begin{theorem}[Cardano 공식]
\label{thm:cardano-formula}
$\charac K\ne2,3$이고
\[
  f(X)=X^3+pX+q
\]
라 하자. 근호들을 포함하는 확장체에서
\[
  u^3=-\frac q2+\sqrt{
  \left(\frac q2\right)^2+\left(\frac p3\right)^3},
\]
\[
  v^3=-\frac q2-\sqrt{
  \left(\frac q2\right)^2+\left(\frac p3\right)^3}
\]
인 $u,v$를
\[
  uv=-\frac p3
\]
가 되도록 택하면 세 근은
\[
  \boxed{x_1=u+v},
\]
\[
  \boxed{x_2=\zeta^2u+\zeta v},
\]
\[
  \boxed{x_3=\zeta u+\zeta^2v}.
\]
\end{theorem}

\begin{proof}
$u=R/3$, $v=S/3$로 두면 명제 \ref{prop:cubic-resolvent-identities}에서 세제곱식과
\[
  uv=\frac{RS}{9}=-\frac p3
\]
을 얻는다. 역 Fourier 변환식에 대입하면 세 근이 나온다.

반대로 표시된 $u,v$가 주어졌다고 하자. $x=u+v$이면
\[
\begin{aligned}
x^3+px+q
&=u^3+v^3+3uv(u+v)+p(u+v)+q\\
&=(-q)+(3uv+p)x+q=0.
\end{aligned}
\]
$u,v$를 $\zeta^2u,\zeta v$ 또는 $\zeta u,\zeta^2v$로 바꾸어도 세제곱과 곱 조건이 유지되므로 나머지 두 근도 얻는다.
\end{proof}

\begin{pitfall}[세제곱근의 독립 선택]
두 세제곱근을 임의로 독립 선택하면 안 된다. 한쪽에 $\zeta$를 곱하면 다른 쪽에는 $\zeta^2$를 곱해야
\[
  uv=-\frac p3
\]
가 유지된다. 이 부조건이 없으면 Cardano 식으로 만든 수가 원래 삼차식의 근이 아닐 수 있다.
\end{pitfall}

\subsection{갈루아군에 따른 단순화}

\begin{itemize}
  \item $\Delta$가 $K$의 제곱이고 $f$가 기약이면 갈루아군은 $A_3\cong C_3$이다. 판별식의 제곱근을 새로 첨가할 필요가 없고, 원시 세제곱근을 준비한 뒤 순환 삼차 근호 단계만 남는다.
  \item $\Delta$가 제곱이 아니고 $f$가 기약이면 갈루아군은 $S_3$이다. 먼저 $\sqrt\Delta$를 첨가한 뒤 세제곱근을 첨가한다.
  \item $f$가 이미 일차인수를 가지면 갈루아군은 더 작은 부분군이고, 한 근을 찾은 뒤 이차식을 풀면 된다.
\end{itemize}

\begin{example}
\[
  X^3-3X-2=(X-2)(X+1)^2
\]
는 분리다항식이 아니므로 위 갈루아군 논의의 기본 가정에서 벗어난다. Cardano 표현은 중근을 포함해 형식적으로 작동할 수 있지만, 분리체의 갈루아군과 근의 세 순열을 그대로 적용하면 안 된다.
\end{example}

\begin{example}
\[
  X^3-3X+1
\]
의 판별식은
\[
  \Delta=-4(-3)^3-27=81
\]
로 제곱이다. 이 다항식은 유리근을 갖지 않으므로 기약이고 갈루아군은 $C_3$이다. 따라서 분해체의 차수는 $3$이며, $S_3$의 이차층은 사라진다.
\end{example}

\begin{keyidea}
Cardano 공식의 제곱근과 세제곱근은 우연히 나타나는 것이 아니다.
\[
  S_3/A_3\cong C_2
  \quad\leadsto\quad
  \sqrt\Delta,
\]
\[
  A_3\cong C_3
  \quad\leadsto\quad
  \sqrt[3]{\phantom{x}}.
\]
공식의 근호 차수는 갈루아군 정규열의 인자군 차수를 그대로 반영한다.
\end{keyidea}

\begin{checkpoint}
\begin{enumerate}[label=(\alph*)]
  \item $X^3-3X+1$에 Cardano 공식을 적용해 $u^3,v^3$를 구하라.
  \item $X^3+3X-4$의 판별식과 갈루아군을 구하고, 유리근을 이용한 풀이와 Cardano 풀이를 비교하라.
  \item $\zeta\notin K$인 경우에도 공식이 근호해법을 주는 이유를 설명하라.
\end{enumerate}
\end{checkpoint}
